% Paper 1 — national post-marketing surveillance of VA-MENGOC-BC, 2017-2025 (English).
% Target journal: Drug Safety (original research article). Build: python tools/latex_build.py
\documentclass[11pt]{article}
\newcommand{\DocLang}{en}
\input{papers/shared/preamble.tex}
\input{papers/shared/authors.tex}
\input{papers/shared/pending.tex}
\input{papers/p1-pharmacovigilance/shared/flow.tex}
\externaldocument{supplement}
\usepackage{lineno}
\onehalfspacing

\hypersetup{pdftitle={Post-marketing safety of the VA-MENGOC-BC meningococcal B and C vaccine in Cuba, 2017-2025},
  pdfauthor={Richard Alejandro Matos Arderi; Abel Ponce Gonzalez}}

\begin{document}

\begin{titlepage}
\SyntheticNotice
\begin{center}
{\LARGE\bfseries Post-marketing safety of the VA-MENGOC-BC meningococcal B and~C vaccine in Cuba,
\VStudyFirstYear--\VStudyLastYear: a nationwide analysis of spontaneous reports with
disproportionality, latent class and machine-learning methods\par}
\bigskip
\AuthorBlock
\end{center}
\bigskip
\noindent\textbf{Running title:} Post-marketing safety of VA-MENGOC-BC in Cuba\par
\noindent\textbf{Keywords:} adverse events following immunization; pharmacovigilance;
disproportionality analysis; meningococcal vaccines; outer membrane vesicle; latent class
analysis; machine learning; Cuba\par
\bigskip
\noindent\fbox{\parbox{\dimexpr\linewidth-2\fboxsep-2\fboxrule}{%
\textbf{Key Points}
\begin{itemize}[leftmargin=*,itemsep=2pt,topsep=2pt]
\item Across \PONReports{} deduplicated reports of adverse events following immunization received
  by the Cuban national surveillance system in \VStudyFirstYear--\VStudyLastYear, \PONTarget{}
  (\POPctTarget) involved VA-MENGOC-BC, with a pooled reporting rate of \POPooledRate{} per
  \num{100000} doses, \PORateVsExpected{} the programme's reference rate for all childhood vaccines.
\item Under a prespecified Bayesian criterion, \POSignalsPrimaryAllOtherVaccines{} of
  \PONEventsScreened{} screened event types were disproportionately reported with
  VA-MENGOC-BC (\POSignalList); \PONSignalsRobustPhrase{} persisted under all
  \PONDesigns{} comparator designs.
\item Co-reported events formed \POLCABestK{} latent phenotypes; hospitalisation was predicted
  with \PODiscriminationWord{} discrimination in later years, so report-based models should
  describe rather than triage.
\end{itemize}}}
\end{titlepage}

\linenumbers

\section*{Abstract}
\noindent\textbf{Introduction:} VA-MENGOC-BC, an outer membrane vesicle vaccine against
serogroup~B and~C meningococcal disease, has been part of the Cuban childhood schedule since
1991, but its post-marketing safety has not been analysed nationally with contemporary
pharmacovigilance methods.\par
\noindent\textbf{Objective:} To describe reporting rates, disproportionality and the structure of
co-reported adverse events following immunization (AEFI) with VA-MENGOC-BC in
\VStudyFirstYear--\VStudyLastYear, and the correlates of hospitalisation.\par
\noindent\textbf{Methods:} Disproportionality analysis of individual case safety reports in the
Cuban national AEFI surveillance database (institutional access), covering
\VStudyFirstYear--\VStudyLastYear{} and all provinces. Reports from \VNFiles{} annual files were
harmonised, pseudonymised and deduplicated. Reporting rates used doses administered as
denominator. VA-MENGOC-BC was compared with all other vaccines using the reporting odds ratio,
proportional reporting ratio, Bayesian information component (IC) and empirical Bayes geometric
mean, with a prespecified signal criterion (lower 95\% credibility bound IC\textsubscript{025}${}>0$
with at least \PODispMinReports{} reports) and three sensitivity comparators; no case-by-case
analysis was performed. Latent class analysis summarised co-reported events; Firth logistic
regression and gradient boosting with temporal validation described hospitalisation. Reporting
followed READUS-PV.\par
\noindent\textbf{Results:} Of \PONReports{} reports, \PONTarget{} (\POPctTarget) involved
VA-MENGOC-BC. The pooled reporting rate was \POPooledRate{} per \num{100000} doses
(\CI{} \POPooledRateCI) and \POTrendWord{} over time (rate ratio per year \POTrendRR; \CI{}
\POTrendRRCI). \POSignalsPrimaryAllOtherVaccines{} of \PONEventsScreened{} event types met the
signal criterion (\POSignalList), \PONSignalsRobust{} of them under every comparator. A
\POLCABestK-class model described co-reporting (bootstrap adjusted Rand index \POLCAAri). The
best model for hospitalisation (\POBestModel) reached an area under the ROC curve of
\POBestAuroc{} in the test years.\par
\noindent\textbf{Conclusions:} This nationwide, reproducible analysis describes the post-marketing
reporting profile of VA-MENGOC-BC. Disproportionality generates hypotheses rather than measuring
risk: signals that persist across comparators warrant clinical review and causality assessment,
and the disclosure-controlled framework supports continuous annual surveillance.

\clearpage
\section{Introduction}\label{sec:intro}

Invasive meningococcal disease remains a cause of death and lifelong sequelae worldwide, with more
than half a million cases and about \num{85000} deaths estimated every year
\cite{SierraGonzalez2019}. In Cuba, the incidence reached 5.9 cases per \num{100000} population in
1980, about 80\% of them caused by serogroup~B, and meningococcal disease was declared the
country's main public health problem \cite{SierraGonzalez2019}. VA-MENGOC-BC, developed at the
Finlay Institute, combines outer membrane vesicles (OMV, proteoliposomes) of a serogroup~B strain
with serogroup~C capsular polysaccharide adsorbed onto aluminium hydroxide
\cite{Sotolongo2007,Uli2006,Masforrol2017}. Its efficacy of 83\% was shown in a randomised,
double-blind, placebo-controlled trial \cite{Sierra1991}; after mass vaccination the vaccine
entered the national immunization programme, and the cumulative impact on serogroup~B disease
exceeded 95\% \cite{Rodriguez1999,SierraGonzalez2019}. Together with the Norwegian and New Zealand
experiences, the Cuban programme is one of the few in which wild-type OMV vaccines controlled
serogroup~B epidemics \cite{Holst2009,Holst2013}.

OMV vaccines have a characteristic reactogenicity profile. Local reactions and fever are frequent
in young children, whereas serious and neurological adverse events are very rare
\cite{Nokleby2007,Holst2009}. The pyrogenicity of OMV has been linked to their membrane-bound
endotoxin, both in experimental models of the four-component serogroup~B vaccine
\cite{Sheerin2019} and in formulation studies, where adsorption onto aluminium hydroxide reduced
pyrogenicity \cite{Rosenqvist1998}. Because VA-MENGOC-BC is administered to infants alongside the
whole-cell pertussis pentavalent vaccine, which is itself the most reactogenic product of the
Cuban schedule \cite{Galindo2012,CruzRodriguez2021}, the attribution of fever and local reactions
to a single product is not straightforward.

Cuba established a national surveillance system for adverse events following immunization (AEFI)
in 1999 \cite{Galindo2012,Reed2007}. Its first decade of data (1999--2008) yielded an overall
rate of 57.8 AEFI per \num{100000} doses \cite{Galindo2012}, and the programme uses 50 per \num{100000}
doses as the expected national rate \cite{CruzRodriguez2021}. Since then, published analyses have
been descriptive and provincial \cite{CruzRodriguez2021}, and no study has examined VA-MENGOC-BC
with the methods now standard in pharmacovigilance: disproportionality analysis with
comparator-sensitivity designs \cite{Evans2001,vanPuijenbroek2002,Bate1998,Noren2013,DuMouchel1999},
explicit handling of masking and co-administration \cite{Maignen2014,Salvo2013}, and transparent
reporting \cite{Fusaroli2024Statement,Fusaroli2024Elaboration}. Machine-learning approaches have
also entered pharmacovigilance, but their added value depends on temporal validation and
calibration rather than on apparent accuracy \cite{Salas2022,Collins2024,VanCalster2019}.

We therefore analysed the national spontaneous reports of \VStudyFirstYear--\VStudyLastYear{}
with four objectives: (i) to
estimate AEFI reporting rates for VA-MENGOC-BC and their trend; (ii) to screen for
disproportionate reporting of specific events against all other vaccines, testing the robustness
of any signal to the choice of comparator; (iii) to describe the latent structure of co-reported
events; and (iv) to describe the correlates of hospitalisation among reports.

\section{Methods}\label{sec:methods}

\subsection{Design, setting and reporting}\label{sec:design}
This was a disproportionality analysis of individual case safety reports, complemented by
reporting-rate, latent class and prediction analyses, in the national AEFI surveillance database of
Cuba. The database receives case investigation forms from vaccination sites in all provinces through
municipal and provincial epidemiology units \cite{Galindo2012,WHO2016AEFI}; its custodian is
\VCustodian{} and the data were extracted on \VExtractionDate. The report follows the READUS-PV guideline for
disproportionality analyses \cite{Fusaroli2024Statement,Fusaroli2024Elaboration} (checklist in
the Supplementary Material), the RECORD-PE statement for routinely collected data
\cite{Langan2018} and, for the hospitalisation models, TRIPOD+AI \cite{Collins2024}.

\subsection{Data sources and processing}\label{sec:data}
The study used \VNFiles{} annual files of AEFI reports (\VStudyFirstYear--\VStudyLastYear) recording,
for each report, demographics, vaccines, doses, lots and manufacturers, dates of
vaccination, onset and notification, \PONEventsScreened{} predefined event indicators, family history,
hospitalisation and outcome. Annual doses of VA-MENGOC-BC administered came from the national
immunization programme series, and the annual incidence of meningococcal disease
(\POITSFirstYear--\POITSLastYear) from national statistics. The database recorded
\PONVaccines{} different vaccines; vaccines were harmonised to a versioned dictionary of
national product names and events were coded with the predefined indicators of the national
form (no MedDRA coding is used by the system). All processing ran on an institutional workstation with an open-source
pipeline (Python 3.11 \cite{Harris2020,Virtanen2020,McKinney2010,Seabold2010,Pedregosa2011,Hunter2007};
version \VCodeVersion, random seed \VSeed) that records the provenance of every run
\cite{Sandve2013,MHRA2018}.

Because the structure of the files changed across years, column names were mapped to a
versioned schema with alias matching and schema-drift reports, and every value was parsed with
an explicit status (valid, recovered, invalid, missing) to support data-quality assessment
following the harmonised framework of Kahn et al. \cite{Kahn2016}. \VFormSentence. Names and addresses were
replaced at ingestion by keyed hashes (HMAC-SHA-256) with a secret key stored outside the
repository \cite{ENISA2019}; dates of birth were converted to age; free-text fields were
processed locally and never released, in accordance with the Cuban data-protection law
\cite{Ley149}. Exact duplicate records (\VExactDuplicates) were removed, and records sharing the
pseudonymised person, vaccination date, vaccines and lots were merged (\VKeyDuplicatesMerged{}
rows), leaving \VReportsDedup{} reports. The analytic year was the year of vaccination; the study
window retained \VReportsInWindow{} reports (\VReportsOutsideWindow{} outside the window were
excluded; \figref{fig:flow}).

\begin{figure}[tbp]
\centering
\ReportFlow
\caption{Flow of AEFI reports from the annual files to the analysis population. Counts come from
the disclosure-controlled release.}
\label{fig:flow}
\end{figure}

\subsection{Exposure and event definitions}\label{sec:events}
A report was attributed to VA-MENGOC-BC when the vaccine was among those administered, including
co-administration with other vaccines. Events were the \PONEventsScreened{} indicators of the national form, which
follow the WHO classification of common, rare and serious events \cite{WHO2016AEFI,CIOMS2012};
when available, the corresponding Brighton Collaboration definitions are cited for reference
\cite{Marcy2004,Bonhoeffer2004Crying,Bonhoeffer2004Seizure,Ruggeberg2007,Kohl2007,Buettcher2007,Wise2007,Gidudu2008,Sejvar2007,Tapiainen2007},
but reports were not re-adjudicated against them \cite{Kohl2008}. Causality assessment by the
national committee \cite{WHO2018Causality} was not used as an inclusion criterion, and no
case-by-case review of the reports was performed in this study. Before the analysis, fever of at
least 39 and 40\,\textdegree{}C, severe local reactions and sterile abscess were classified as
expected reactions to aluminium-adsorbed OMV vaccines \cite{Nokleby2007,Holst2009,Jefferson2004};
signals for any other event were classified as emerging.

\subsection{Reporting rates}\label{sec:m-rates}
Annual reporting rates were computed per \num{100000} doses of VA-MENGOC-BC with exact Poisson
(Garwood) intervals \cite{Garwood1936}. The linear trend was estimated with a quasi-Poisson model
with log-dose offset (rate ratio per year, dispersion estimated from Pearson's statistic) and
checked with a negative-binomial model. Rates were compared descriptively with the programme's
expected rate of \POExpectedRate{} per \num{100000} doses, which refers to all childhood vaccines
\cite{CruzRodriguez2021,Mahaux2016}.

\subsection{Disproportionality analysis}\label{sec:m-disp}
For each event, a two-by-two table crossed reports with and without VA-MENGOC-BC against
reports with and without the event. We computed the reporting odds ratio (ROR) with Woolf
intervals and Haldane correction when needed \cite{vanPuijenbroek2002,Woolf1955,Haldane1956},
the proportional reporting ratio (PRR) with Yates' $\chi^2$ \cite{Evans2001}, the Bayesian
confidence propagation neural network information component (IC) with its lower 95\% credibility
bound IC\textsubscript{025} \cite{Bate1998,Noren2013}, and the empirical Bayes geometric mean
(EBGM) of the multi-item gamma Poisson shrinker, whose two-gamma prior was estimated by maximum
likelihood over all vaccine--event pairs of the database \cite{DuMouchel1999,Szarfman2002}. The
prespecified primary signal criterion was IC\textsubscript{025}${}>0$ with at least
\PODispMinReports{} reports; the number of methods meeting their conventional thresholds (ROR lower bound ${}>1$;
PRR ${}\geq2$ with $\chi^2\geq4$; EBGM lower bound ${}\geq2$) is reported as a measure of
concordance.

All comparators were vaccines, so the measures express reporting relative to other vaccines in
the same system rather than risk relative to non-vaccination. Four comparator designs were
prespecified: the primary design (VA-MENGOC-BC reports against all other reports); an infant
active-comparator design restricted to reports in children younger than \POInfantMonths{}
months; a design
restricted to single-vaccine reports (VA-MENGOC-BC given alone against other vaccines given
alone); and a design removing the reports that mention the pentavalent vaccine, the dominant
reactogenic product, from the comparator group to assess masking \cite{Maignen2014}. The cumulative IC was
recomputed year by year to show the temporal stability of the estimates. Given the exploratory
nature of signal detection, no adjustment for multiplicity was applied to the screening
\cite{Rothman1990}; signals are hypotheses for clinical review rather than confirmed risks
\cite{Fusaroli2026Pitfalls,Fusaroli2026Causal}.

\subsection{Latent class analysis of co-reported events}\label{sec:m-lca}
Event indicators with a prevalence of at least 0.5\% were modelled with latent class analysis (a
mixture of independent Bernoulli distributions) fitted by the expectation--maximisation algorithm
\cite{Dempster1977} with a Beta(1.5, 1.5) prior on item probabilities and \POLCAStarts{} random
starts per model. Models with \POLCAKMin--\POLCAKMax{} classes were compared with the Bayesian information criterion
\cite{Schwarz1978}; the integrated completed likelihood \cite{Biernacki2000} and entropy were
reported \cite{Nylund2007}. Stability was measured by the median adjusted Rand index
\cite{Hubert1985} between the original classification and \POLCABootstrap{} bootstrap refits, after optimal
label matching \cite{Kuhn1955}. Class membership was compared between VA-MENGOC-BC and other
reports with a $\chi^2$ test.

\subsection{Correlates of hospitalisation}\label{sec:m-hosp}
Hospitalisation was modelled from age, sex, co-administration, region, vaccines and event
indicators. Because rare, serious events separate the outcome, the multivariable logistic model
was fitted by Firth's penalised likelihood \cite{Firth1993,Heinze2002}, after removal of exactly
collinear terms, and its adjusted odds ratios are reported with Wald intervals. A gradient-boosted
tree model (LightGBM) \cite{Ke2017} captured non-linearities and interactions. Both models were
trained on \POTrainYears{} and evaluated on \POTestYears{} (temporal validation), reporting the area under
the ROC and precision--recall curves, the Brier score and the calibration intercept and slope
\cite{VanCalster2019,Collins2024}; logistic predictions used the FLIC intercept correction
\cite{Puhr2017}. Exact TreeSHAP contributions described the boosted model
\cite{Lundberg2020}. The models are explanatory and were not designed for clinical triage.

\subsection{Temporal patterns and epidemiological context}\label{sec:m-time}
Monthly report counts were variance-stabilised with the Anscombe transformation
\cite{Anscombe1948} and segmented with an exact PELT algorithm \cite{Killick2012,Truong2020}
(minimum segment of six months, BIC-type penalty). To place the vaccine in context, the annual
national cases of meningococcal disease (\POITSFirstYear--\POITSLastYear) were analysed as an
interrupted time series with a segmented negative-binomial model, treating \POITSTransition{}
(mass vaccination and inclusion in the programme) as a transition period \cite{Bernal2017}.

\subsection{Disclosure control, reproducibility and ethics}\label{sec:m-sdc}
Only aggregated results left the controlled environment. Non-zero counts below \VMinCell{}
were suppressed, statistics derived from suppressed counts were removed and
complementary suppression protected contingency tables \cite{Hundepool2012}; an automated
verifier checked every released file. Every number in this article is generated from the
released tables, and the pipeline, synthetic test data and release are openly available
\cite{Wilkinson2016}. The analysis used pseudonymised routine surveillance records processed
under the institutional data-governance procedure of the Finlay Vaccine Institute and Law
149/2022 \cite{Ley149}; the institutional approval is detailed in the Declarations.

\section{Results}\label{sec:results}

\subsection{Reports and characteristics}\label{sec:r-reports}
The analysis included \PONReports{} reports, of which \PONTarget{} (\POPctTarget) involved
VA-MENGOC-BC and \PONTargetOnly{} of these had no co-administered vaccine
(\tabref{tab:characteristics}). The median delay from vaccination to notification was
\PODelayMedian{} days.

\begin{table}[tbp]
\centering\small
\caption{Characteristics of AEFI reports involving VA-MENGOC-BC and other vaccines. Values are
$n$ (\%) unless stated otherwise; counts below \VMinCell{} are suppressed, and [c] marks
complementary suppression that prevents their recovery by subtraction.}
\label{tab:characteristics}
\POTableCharacteristics
\end{table}

\subsection{Reporting rates}\label{sec:r-rates}
The pooled reporting rate for VA-MENGOC-BC between \PORateFirstYear{} and \PORateLastYear{} was
\POPooledRate{} per \num{100000} doses (\CI{} \POPooledRateCI), \PORateVsExpected{} the programme's
reference rate of \POExpectedRate{} per \num{100000} doses for all childhood vaccines
(\figref{fig:rates}; Supplementary \tabref{tab:S-rates}). The rate \POTrendWord{} over the period (rate ratio
per year \POTrendRR, \CI{} \POTrendRRCI; $p$\POTrendP; quasi-Poisson dispersion
\PODispersion).

\begin{figure}[tbp]
\centering
\includegraphics[width=\textwidth]{p1_rates}
\caption{Reporting over time. (a) Monthly number of AEFI reports for all vaccines and for
VA-MENGOC-BC; vertical lines mark the change points detected by PELT in the series of all
vaccines. (b) Annual reporting rate of VA-MENGOC-BC per \num{100000} doses with exact 95\%
confidence interval (band); the horizontal line is the programme's expected rate for all childhood
vaccines.}
\label{fig:rates}
\end{figure}

\subsection{Disproportionality}\label{sec:r-disp}
Under the primary comparator, \POSignalsPrimaryAllOtherVaccines{} of the \PONEventsScreened{}
screened event types met the prespecified signal criterion (\POSignalList;
\tabref{tab:disproportionality}, \figref{fig:forest}): \PONSignalsExpectedPhrase{} for expected
reactions (\POSignalExpectedList) and \PONSignalsEmergingPhrase{} for events not classified as
expected (\POSignalEmergingList). The number of signals was \POSignalsInfantActiveComparator{} in
the infant active-comparator design, \POSignalsExcludingCoadministration{} in single-vaccine
reports and \POSignalsExcludingPentavalentMasking{} when pentavalent-vaccine reports were removed
from the comparator (\tabref{tab:designs}; Supplementary
Tables~\ref{tab:S-disp-primary}--\ref{tab:S-disp-masking}); \PONSignalsRobustPhrase{} persisted in
all \PONDesigns{} designs (\POSignalRobustList). The cumulative IC of each event is shown year by
year in Supplementary \tabref{tab:S-cumic}. For fever of at least 39\,\textdegree{}C, the ROR
was \POEvFeverThreeNineROR{} (\CI{} \POEvFeverThreeNineRORCI) and IC\textsubscript{025}
\POEvFeverThreeNineICLow; for severe local reactions they were \POEvSevereLocalReactionROR{}
(\CI{} \POEvSevereLocalReactionRORCI) and \POEvSevereLocalReactionICLow.

\begin{table}[tbp]
\centering\small
\caption{Disproportionality of events reported with VA-MENGOC-BC against all other vaccines
(primary design), for events with at least \VMinCell{} reports. $a$, reports with VA-MENGOC-BC and
the event; $E$, expected $a$; IC\textsubscript{025}, lower 95\% credibility bound of the
information component; EB\textsubscript{05}, lower 5th percentile of the empirical Bayes
geometric mean; Methods, number of the four methods meeting their thresholds; Signal,
prespecified criterion (IC\textsubscript{025}${}>0$, $a\geq\PODispMinReports$).}
\label{tab:disproportionality}
\POTableDispMain
\end{table}

\begin{table}[tbp]
\centering\small
\caption{Number of disproportionality signals under the four prespecified comparator designs.}
\label{tab:designs}
\POTableDesigns
\end{table}

\begin{figure}[tbp]
\centering
\includegraphics{p1_forest}
\caption{Reporting odds ratios (95\% CI, logarithmic scale) of events reported with VA-MENGOC-BC
against all other vaccines (primary design) and in the infant active-comparator design. Events
with fewer than \VMinCell{} reports in the primary design are not shown; all four designs are
reported in Supplementary Tables~\ref{tab:S-disp-primary}--\ref{tab:S-disp-masking}.}
\label{fig:forest}
\end{figure}

\subsection{Latent phenotypes of co-reported events}\label{sec:r-lca}
The Bayesian information criterion selected a \POLCABestK-class model (Supplementary
\tabref{tab:S-lca}),
whose classification was stable across bootstrap samples (median adjusted Rand index
\POLCAAri). Class membership \POLCAGroupWord{} between VA-MENGOC-BC reports and other reports
($\chi^2=\POLCAChi$, \POLCADf{} degrees of freedom; $p$\POLCAP; \figref{fig:lca}).

\begin{figure}[tbp]
\centering
\includegraphics[width=\textwidth]{p1_lca}
\caption{Latent classes of co-reported events. (a) Probability of each event within each latent
class (values of at least 0.2 are printed). (b) Share of reports assigned to each class among
reports with VA-MENGOC-BC and with other vaccines.}
\label{fig:lca}
\end{figure}

\subsection{Correlates of hospitalisation}\label{sec:r-hosp}
The training period included \PONTrain{} reports (\POEventsTrain{} hospitalisations) and the test
period \PONTest{} reports (\POEventsTest{} hospitalisations). In the test years, the
\POBestModel{} model had the higher discrimination (area under the ROC curve \POBestAuroc;
\tabref{tab:models}), which we graded as \PODiscriminationWord. A calibration slope below one
indicates predictions that are too extreme in the test years, and a slope above one predictions
that are too moderate. The adjusted odds ratios of the Firth model are given in Supplementary
\tabref{tab:S-logistic} and the TreeSHAP importance in \figref{fig:shap}.

\begin{table}[tbp]
\centering\small
\caption{Temporal validation of the hospitalisation models (trained on \POTrainYears, tested on
\POTestYears). AUROC and AUPRC, areas under the receiver operating characteristic and
precision--recall curves.}
\label{tab:models}
\POTableModels
\end{table}

\begin{figure}[tbp]
\centering
\includegraphics{p1_shap}
\caption{Mean absolute TreeSHAP contribution of the twelve most influential predictors to the
predicted log-odds of hospitalisation in the test years (gradient-boosted model). Region
contrasts are relative to the western region.}
\label{fig:shap}
\end{figure}

\subsection{Temporal patterns and context}\label{sec:r-time}
The PELT segmentation identified \PONBreaksAllPhrase{} in the monthly number of reports
(\POBreaksAll). In the national series of meningococcal disease, the segmented model showed
\POITSWord{} after the \POITSTransition{} transition (level rate ratio \POITSLevelRR, \CI{}
\POITSLevelRRCI); the mean incidence was \POITSRatePre{} per \num{100000} population before the
campaign and \POITSRatePost{} afterwards (Supplementary \figref{fig:S-its}).

\section{Discussion}\label{sec:discussion}

\subsection{Principal findings}\label{sec:d-principal}
In \VStudyFirstYear--\VStudyLastYear, VA-MENGOC-BC accounted for \POPctTarget{} of AEFI
reports, with a pooled reporting rate of \POPooledRate{} per \num{100000} doses that
\POTrendWord{} over the period. Disproportionality analysis flagged \POSignalList{} under the
primary comparator, and \POSignalRobustList{} under every comparator. Co-reported events organised
into \POLCABestK{} phenotypes, and the information in the reports predicted hospitalisation with
\PODiscriminationWord{} discrimination.

\subsection{Interpretation}\label{sec:d-interpretation}
The most frequently reported events with VA-MENGOC-BC were \POTopEvents. Fever and local
reactions are the expected reactions to aluminium-adsorbed OMV vaccines \cite{Nokleby2007,Holst2009},
whose endotoxin content drives pyrogenicity \cite{Sheerin2019,Rosenqvist1998}. Disproportionality in a
vaccine-only database measures how the profile of VA-MENGOC-BC reports differs from that of other
vaccines. Because most comparators are themselves reactogenic, and the pentavalent whole-cell
vaccine is co-administered with VA-MENGOC-BC, a signal that disappears when co-administration or
the pentavalent vaccine is removed is better explained by attribution and masking than by the
product \cite{Maignen2014,Salvo2013}; conversely, signals that persist across all comparators are
those most deserving clinical review. The structured, multi-comparator design is therefore the
main safeguard of this analysis, together with the prespecified criterion and the transparent
reporting of every event screened \cite{Fusaroli2024Elaboration,Fusaroli2026Pitfalls}.
Signals for expected reactions (\POSignalExpectedList) quantify how much more often the known
reactogenicity of an OMV vaccine appears in reports than for other vaccines; signals for emerging
events (\POSignalEmergingList) are the ones that require case-by-case review before any
interpretation.

The reporting rate must be read against the characteristics of the Cuban system: a mandatory form
completed at the vaccination site, active follow-up by epidemiology services and a strong
reporting culture \cite{Galindo2012,Reed2007}. Reporting rates in such systems reflect the
sensitivity of surveillance as much as the frequency of reactions, as the Pinar del Río experience
showed \cite{CruzRodriguez2021}, and
the WHO uses reporting rates as an indicator of surveillance performance \cite{Lei2018}. The
programme's reference rate refers to all vaccines, so the comparison with VA-MENGOC-BC is only a
benchmark \cite{Mahaux2016}.

Latent classes offer a data-driven alternative to single-event screening: they describe
syndromes as they are reported, separate febrile from local or allergic presentations and allow
the comparison of phenotype distributions between vaccines. The hospitalisation models reached
\PODiscriminationWord{} discrimination in later years. \POHospInterpretation, so report-based models
should not be used to triage cases \cite{VanCalster2019,Collins2024}.

The findings apply to a mature, mandatory surveillance system with high vaccination coverage and
to an aluminium-adsorbed wild-type OMV vaccine given with a whole-cell pentavalent vaccine; their
transfer to other settings or OMV formulations requires the same comparator-sensitivity approach.
Two further designs would test the signals: a self-controlled case series in linked immunisation
and hospital records, which removes time-invariant confounding, and the analysis of the quality
attributes of the administered lots, which addresses manufacturing variability.

\subsection{Strengths and limitations}\label{sec:d-limitations}
The study covers every report received nationally in \VStudyFirstYear--\VStudyLastYear, harmonises files whose
structure changed over time, prespecifies its signal criterion and comparator designs, and
generates every number from a verified, disclosure-controlled release, so that the analysis can be
reproduced and extended each year.

Its limitations are those of spontaneous reporting, reinforced by specific features of the data.
First, under-reporting and stimulated reporting affect the numerators \cite{Hazell2006,Pariente2007};
doses administered, the only denominator available, were not stratified by age, dose or province.
Second, the event indicators were not validated against Brighton definitions
\cite{Kohl2008}, and free-text descriptions were excluded from the analysis to protect
confidentiality. Third, reports attributed to VA-MENGOC-BC include co-administered vaccines, and
disproportionality cannot separate their contributions; the sensitivity designs address but do
not remove this limitation. Fourth, all comparators are vaccines, so disproportionality does not
measure risk relative to non-vaccination, and a signal is a hypothesis for causality assessment
\cite{WHO2018Causality,Fusaroli2026Causal}. Fifth, the interrupted time series of national
incidence is ecological and serves only as context. Finally, small counts were suppressed in the
release, so rare serious events can be described only qualitatively.

\subsection{Implications}\label{sec:d-implications}
The pipeline developed for this study can run each year on new files, turning a retrospective
analysis into continuous, auditable surveillance. Signals that persist across comparators should
be reviewed with the clinical records by the national causality committee, and linking reports to
the quality attributes of the administered lots offers a way to study whether manufacturing
variability contributes to reactogenicity.

\section{Conclusions}\label{sec:conclusions}
This nationwide analysis describes the post-marketing reporting profile of VA-MENGOC-BC with
methods that are transparent about their assumptions and reproducible from a public,
disclosure-controlled release. Signals that persist across comparator designs define the
priorities for clinical review, and the framework provides a basis for continuous safety
surveillance of the Cuban immunization programme.

\section*{Declarations}
\input{papers/p1-pharmacovigilance/shared/declarations_en.tex}

\bibliographystyle{vancouver}
\bibliography{papers/shared/bib/pubmed,papers/shared/bib/manual}

\end{document}
